\documentclass[10pt,a4paper]{amsart}
\usepackage[margin=19mm]{geometry}
\usepackage{amsmath,amssymb,mathtools}
\usepackage[scr=rsfs,cal=euler]{mathalfa}
\usepackage{xfrac}
\usepackage[unicode,hidelinks,bookmarks=false]{hyperref}
\newcommand{\ie}{i.e.\ }
\newcommand{\cf}{cf.\ }
\newcommand{\resp}{respectively\ }
\newcommand{\Subsection}{\subsection}
\providecommand{\nobreakdash}{\nobreak\hbox{-}}
\setlength{\emergencystretch}{2em}
\multlinegap0pt
\usepackage{amssymb}
\usepackage[shortlabels]{enumitem}
\newtheorem{theorem}{Theorem}
\title{The quaternionic moment problem}
\author{R. Ben Taher, K. Schm\"udgen, E.H. Zerouali}
\date{Source: arXiv:2608.00188v1 (CC BY 4.0)}
\begin{document}
\maketitle

\noindent\emph{Source and licence.} The quaternionic moment problem. Version arXiv:2608.00188v1 (CC BY 4.0), distributed by arXiv under the Creative Commons Attribution 4.0 International licence, http://creativecommons.org/licenses/by/4.0/, as recorded in the arXiv licence metadata for this exact version and shown on the abstract page https://arxiv.org/abs/2608.00188v1. Full licence notice: \texttt{licenses/arxiv-2608.00188-CC-BY-4.0.txt}.\par\smallskip

\noindent\textit{Editorial scope.} Counted unit: the article's theorem realpartH\_L (source0.tex lines 898--914). The article's complete original proof of it is delivered with it (source0.tex lines 916--1075, one \texttt{proof} environment of 160 lines). No mathematics is written, repaired or re-derived: the statement and the proof are the article's own lines, in the article's own notation, and no retained line is substituted or re-flowed.  Retained uncounted as the source-local context the proof needs: lines 328--331.  Nothing was omitted from any retained span, so the package carries no omission record.  No retained line contains a citation, so the wrapper carries no bibliography and no bibliography entry.  No retained line contains a figure environment, an image-inclusion command or drawing code, so no image is delivered and the package declares no asset dependency.  Editorial formatting only, in addition: an AMS \texttt{amsart} preamble that restates the article's own declarations and macros used by the retained lines, namely the environments \texttt{theorem} on separate counters, together with the standard packages those lines need.  The retained environments print in this document as Theorem 1 in order of appearance, whereas the article prints the same items with the numbering of its own full document.  The complete native source of the article as distributed in the arXiv e-print for this exact version is bundled unchanged as \texttt{sources/206522/source0.tex}.
\begin{align}
\label{EherRe}
E_{\mathrm{her}}=\{\, \mathrm{Re}f: \, f\in E\}.
\end{align}

\begin{theorem}%
\label{realpartH_L}
%
\begin{align}
\label{herpart}
\mathbb{H}^{0}_{\mathbb{R}}[q^{*}, q]_{\mathrm{her}} =& ~\mathbb{R}[x_{0},x_{1}^{2}+x_{2}^{2}+x_{3}^{2}],
\\
\mathbb{H}^{0}[q^{*}, q]_{\mathrm{her}} =& ~\mathbb{R}[x_{0},x_{1}^{2}+x_{2}^{2}+x_{3}^{2}]
+ x_{1}\, \mathbb{R}[x_{0},x_{1}^{2}+x_{2}^{2}+x_{3}^{2}]
\label{herpart1}
\\
&+x_{2}\, \mathbb{R}[x_{0},x_{1}^{2}+x_{2}^{2}+x_{3}^{2}]+ x_{3}\,
\mathbb{R}[x_{0},x_{1}^{2}+x_{2}^{2}+x_{3}^{2}].
\nonumber
\end{align}
%
\end{theorem}

\begin{proof}
We write $q=x_{0}+x_{I}$, where
$x_{I}=\mathrm{i}x_{1}+\mathrm{j}x_{2}+\mathrm{k}x_{3}$. Then
$q^{*}=x_{0}-x_{I}$.

Since $x_{0}$ and $x_{I}$ commute, the binomial theorem applies and we
obtain
%
\begin{align*}
(q^{*})^{k}q^{n} &= (x_{0}-x_{I})^{k}(x_{0}+x_{I})^{n}
\\
&=\left (\sum _{j=0}^{k}\binom{k}{j} x_{0}^{k-j}(-x_{I})^{j}\right )
\left (\sum _{\ell =0}^{n}\binom{n}{\ell} x_{0}^{n-\ell}x_{I}^{\ell}
\right )
\\
&= \sum _{j=0}^{k} \sum _{\ell =0}^{n} \binom{k}{j}\binom{n}{\ell} (-1)^{j}
x_{0}^{k+n-j-\ell} x_{I}^{j+\ell}.
\end{align*}

For $r\in \mathbb{N}$, we have
%
\begin{align*}
x_{I}^{2r}&=(-1)^{r}\|x_{I}\|^{2r} = (-1)^{r}(x_{1}^{2}+x_{2}^{2}+x_{3}^{2})^{r},
\\
x_{I}^{2r+1}&=(-1)^{r}\|x_{I}\|^{2r} x_{I} = (-1)^{r}(x_{1}^{2}+x_{2}^{2}+x_{3}^{2})^{r}
x_{I}.
\end{align*}
%
Hence, $\mathrm{Re}((q^{*})^{k}q^{n})$ is a real combination of terms
$x_{0}^{k+n-j-\ell}x_{I}^{j+\ell}$ with $j+\ell $ even, and
$\mathrm{Im}((q^{*})^{k}q^{n})$ is a real combination of terms
$x_{0}^{k+n-j-\ell}x_{I}^{j+\ell}$ with $j+\ell $ odd. Therefore,
%
\begin{align}
\label{re}
\mathrm{Re}((q^{*})^{k}q^{n}) & \in \mathbb{R}[x_{0},x_{I}^{2}]
\subseteq \mathbb{R}[x_{0},x_{1}^{2}+x_{2}^{2}+x_{3}^{2}],
\\
\mathrm{Im}((q^{*})^{k}q^{n}) & \in x_{I} \, \mathbb{R}[x_{0},x_{I}^{2}]
\subseteq x_{I} \, \mathbb{R}[x_{0},x_{1}^{2}+x_{2}^{2}+x_{3}^{2}].
\label{im}
\end{align}
%
In particular, from (\ref{EherRe}) and (\ref{re}) we obtain
%
\begin{align}
\label{hrinclusion}
\mathbb{H}^{0}_{\mathbb R}[q^{*}, q]_{\mathrm{her}} \subseteq
\mathbb{R}[x_{0},x_{1}^{2}+x_{2}^{2}+x_{3}^{2}].
\end{align}

Now we prove the converse inclusion in (\ref{hrinclusion}). Let
$p,s\in \mathbb{N}$. Then
%
\begin{align}
x_{0}^{p} & = \frac{1}{2^{p}}\, (q+q^{*})^{p} = \frac{1}{2^{p}} \sum _{k=0}^{p}
\binom{p}{k} q^{k} (q^{*})^{p-k},
\label{x0}
\\
(x_{1}^{2}+x_{2}^{2}+x_{3}^{2})^{s} & = (q^{*}q - x_{0}^{2})^{s} =
\sum _{j=0}^{s} \binom{s}{j} (-1)^{s-j} (q^{*})^{j} q^{j} (x_{0}^{2})^{s-j}.
\label{xsum}
\end{align}
%
(Note that the binomial theorem applies, because $q,q^{*}$ and
$q^{*}q, x_{0}^{2}$ commute.) By (\ref{x0}), $(x_{0}^{2})^{s-j}$ is a real
combination of terms $(q^{*})^{k}q^{n}$. Therefore, since
$\mathbb{H}^{0}_{\mathbb R}[q^{*}, q]$ is a commutative algebra, it follows
from (\ref{x0}) and (\ref{xsum}) that the product
$x_{0}^{p}(x_{1}^{2}+x_{2}^{2}+x_{3}^{2})^{s}$ is a real combination of
terms $(q^{*})^{k}q^{n}$, so it is contained in
$ \mathbb{H}^{0}_{\mathbb{R}}[q^{*}, q]$. Because it is real-valued on
$\mathbb{H}$,
$x_{0}^{p}(x_{1}^{2}+x_{2}^{2}+x_{3}^{2})^{s} \in \mathbb{H}^{0}_{
\mathbb R}[q^{*}, q]_{\mathrm{her}}$. Hence
$ \mathbb{R}[x_{0},x_{1}^{2}+x_{2}^{2}+x_{3}^{2}]\subseteq \mathbb{H}^{0}_{
\mathbb R}[q^{*}, q]_{\mathrm{her}} $. Combined with (\ref{hrinclusion}),
this proves (\ref{herpart}).

Now we turn to the proof of (\ref{herpart1}). For short, we denote the
set at the right-hand side of (\ref{herpart1}) by $\mathcal{R}$. Then the
assertion (\ref{herpart1}) means that
$ \mathbb{H}^{0}_{\mathbb R}[q^{*}, q]_{\mathrm{her}}=\mathcal{R}$.

First we show that
$ \mathbb{H}^{0}[q^{*}, q]_{\mathrm{her}}\subseteq \mathcal{R}$. By (\ref{EherRe})
it suffices to prove that
%
\begin{align}
\label{rekn}
\mathrm{Re} (\tau [ (q^{*})^{k}q^{n} ]\sigma ) \in \mathcal{R}\quad
\text{ for all} ~~~ k,n\in \mathbb{N}~~~\text{and}~~\tau , \sigma \in
\{1,i,j,k\}.
\end{align}
%
Since
$\mathrm{Re}[(q^{*})^{k}q^{n}] \in \mathbb{R}[x_{0},x_{1}^{2}+x_{2}^{2}+x_{3}^{2}]$
as noted above, we conclude that
%
\begin{align}
\label{reqkn}
\mathrm{Re} (\tau \, \mathrm{Re} [ (q^{*})^{k}q^{n} ]\,\sigma )=
\mathrm{Re} [ (q^{*})^{k}q^{n} ]\, \cdot \, \mathrm{Re} (\tau \sigma )
\in \mathcal{R}.
\end{align}
%
For the imaginary part we have
$\mathrm{Im}[(q^{*})^{k}q^{n}] \in x_{I} \, \mathbb{R}[x_{0},x_{1}^{2}+x_{2}^{2}+x_{3}^{2}]$
by (\ref{im}). Hence
%
\begin{align}
\label{imqkn1}
\mathrm{Re} (\tau \, \mathrm{Im} \, [ (q^{*})^{k}q^{n} ]\, \sigma )=
\mathbb{R}[x_{0},x_{1}^{2}+x_{2}^{2}+x_{3}^{2}]\, \cdot \,
\mathrm{Re} (\tau x_{I} \sigma ).
\end{align}
%
Since
$\tau x_{I} \sigma =\tau i \sigma x_{1}+ \tau j\sigma x_{2}+ \tau k
\sigma x_{3}$, $\mathrm{Re} (\tau x_{I} \sigma )$ is a real linear combination
of $x_{1},x_{2},x_{3}$. Thus it follows from (\ref{imqkn1}) that
%
\begin{align}
\label{imqkn}
\mathrm{Re} (\tau \, \mathrm{Im} [ (q^{*})^{k}q^{n} ]\,\sigma )\in
\mathcal{R}.
\end{align}
%
Clearly, (\ref{reqkn}) and (\ref{imqkn}) imply (\ref{rekn}) which proves
that $ \mathbb{H}^{0}[q^{*}, q]_{\mathrm{her}}\subseteq \mathcal{R}$.

Now we verify the opposite inclusion
$\mathcal{R}\subseteq \mathbb{H}^{0}[q^{*}, q]_{\mathrm{her}}$. Suppose
that $f\in \mathbb{R}[x_{0},x_{1}^{2}+x_{2}^{2}+x_{3}^{2}]$. Set
$g:=\frac{1}{2}(q^{*}\mathrm{i}-\mathrm{i}q) f$. Since
$f\in \mathbb{H}_{\mathbb{R}}[q^{*}, q]_{\mathrm{her}}$ by (\ref{herpart}),
we have $g\in \mathbb{H}^{0}[q^{*}, q]$ and
$\mathrm{Re} ~ g= x_{1} f$. This shows that
%
\begin{equation*}
x_{1} \mathbb{R}[x_{0},x_{1}^{2}+x_{2}^{2}+x_{3}^{2}] \subseteq
\mathbb{H}^{0}[q^{*}, q]_{\mathrm{her}}.
\end{equation*}
%
Replacing $x_{1}= \frac{1}{2}(q^{*}\mathrm{i}-\mathrm{i}q)$ by
$x_{2}= \frac{1}{2}(q^{*}\mathrm{j}- \mathrm{j}q) $ and
$x_{3}= \frac{1}{2}(q^{*}\mathrm{k}-\mathrm{k}q) $, respectively, the same
reasoning yields
%
\begin{equation*}
x_{2} \mathbb{R}[x_{0},x_{1}^{2}+x_{2}^{2}+x_{3}^{2}] \subseteq
\mathbb{H}^{0}[q^{*}, q]_{\mathrm{her}}, ~~~ x_{3} \mathbb{R}[x_{0},x_{1}^{2}+x_{2}^{2}+x_{3}^{2}]
\subseteq \mathbb{H}^{0}[q^{*}, q]_{\mathrm{her}}.
\end{equation*}
%
Therefore,
$\mathcal{R}\subseteq \mathbb{H}^{0}[q^{*}, q]_{\mathrm{her}}$. Thus, together
with the preceding paragraph we have shown that equation (\ref{herpart1})
holds. This completes the proof of Theorem~\ref{realpartH_L}.
\end{proof}

\end{document}
